% Chapter 2 -- Completing Graphs (key: graphon)
\casechapter{Completing Graphs}{From the regularity lemma to graph limits}{Completion}{ch:graphon}
\begin{refsection}[chapters/02-graphon.bib]

% Chapter-specific notation.
\newcommand{\graphoncut}[1]{\norm{#1}_{\square}}
\newcommand{\graphondelta}{\delta_{\square}}
\newcommand{\graphonhom}{\operatorname{hom}}
\newcommand{\graphonspace}{\widetilde{\cW}_0}

\begin{chaptersummary}
Extremal graph theory asks how the density of one small graph in a large graph constrains the density of another. For decades its theorems were statements about sequences of ever larger graphs, proved with $\eps$'s, regular partitions, and towers of exponentials, and its optimal configurations were sequences rather than objects. Lov\'asz and Szegedy completed the space of graphs: a graph becomes a black-and-white picture on the unit square, and the limits of such pictures are the symmetric measurable functions $W\colon[0,1]^2\to[0,1]$, under a metric in which subgraph densities are continuous and the space is compact. In the completion, quasirandomness is convergence to a constant function, the counting lemma loses its hypotheses, and extremal problems acquire extremal objects. The price is a compactness theorem, and the bounds it does not give.
\end{chaptersummary}

\section{Before: extremal graph theory with epsilons}\label{sec:graphon:before}

Dense graph theory measures a large graph by how often small graphs occur in it. All graphs are finite and simple, and $v(\cdot)$ and $e(\cdot)$ count vertices and edges. A \term{homomorphism} from $F$ to $G$ is a map $V(F)\to V(G)$ that sends edges to edges. With $\graphonhom(F,G)$ their number, the \term{homomorphism density}\index{homomorphism density} $t(F,G)\defeq\graphonhom(F,G)/v(G)^{v(F)}$ is the probability that a uniformly random map is a homomorphism.\sidenote{Counting homomorphisms rather than copies, as Lov\'asz and Szegedy do, makes $t(F,G)$ a probability and makes it multiplicative: the density of a disjoint union $F_1\sqcup F_2$ is $t(F_1,G)\,t(F_2,G)$.} If $G$ has $n$ vertices and $T$ triangles, then $t(K_2,G)=2e(G)/n^2$ and $t(K_3,G)=6T/n^3$. Maps that are not injective are a fraction at most $\binom{v(F)}{2}/n$ of all maps, so for large $G$ counting homomorphisms or copies makes no difference.

\subsection{Extremal problems at a given edge density}

Extremal graph theory asks which combinations of densities are possible. Its first theorem is due to \textcite{Mantel1907}:\sidenote{It appeared as the solution of Problem~28 in the Dutch problem journal \emph{Wiskundige Opgaven}. Tur\'an extended it in 1941 from triangles to all cliques.} a triangle-free graph on $n$ vertices has at most $n^2/4$ edges, so $t(K_3,G)=0$ forces $t(K_2,G)\le1/2$. Above edge density $1/2$ triangles must appear, and Goodman found how many.\sidenote{Goodman's paper is about parties: how many triples of guests must be mutual acquaintances or mutual strangers? Asymptotically, at least a quarter of all triples.}\index{Goodman's bound}

\begin{theorem}[\textcite{Goodman1959}]\label{thm:graphon:goodman-finite}
Every graph $G$ satisfies $t(K_3,G)\ge t(K_2,G)\bigl(2t(K_2,G)-1\bigr)$.
\end{theorem}

\begin{proof}
Let $G$ have $n$ vertices, $m$ edges, $T$ triangles, and degrees $d(v)$. For an edge $uv$, the neighbourhoods $N(u)$ and $N(v)$ lie in a set of $n$ vertices, so $\abs{N(u)\cap N(v)}\ge d(u)+d(v)-n$. Each edge $uv$ lies in exactly $\abs{N(u)\cap N(v)}$ triangles and each triangle has three edges, so summing over the edges gives
\[
  3T\;\ge\sum_{uv\in E(G)}\bigl(d(u)+d(v)\bigr)-mn\;=\;\sum_{v}d(v)^2-mn\;\ge\;\frac{4m^2}{n}-mn .
\]
The equality holds because each vertex $v$ lies in $d(v)$ edges, and the last step is the Cauchy--Schwarz inequality $\sum_v d(v)^2\ge\bigl(\sum_v d(v)\bigr)^2/n$. Since $t(K_3,G)=6T/n^3$ and $p\defeq t(K_2,G)=2m/n^2$, the inequality $3T\ge m(4m-n^2)/n$ is $t(K_3,G)\ge p(2p-1)$.
\end{proof}

The bound is attained at $p=1-1/r$ by complete $r$-partite graphs with equal parts (Exercise~\ref{ex:graphon:region}), but not in between, where the true minimum was found only in 2008 (\cref{thm:graphon:razborov}). Even to state that minimum takes a limit: it is the least possible limit of $t(K_3,G_n)$ over sequences of graphs whose edge densities tend to $p$. The extremal objects are sequences.

Sometimes no sequence can be singled out. The Cauchy--Schwarz inequality, applied to codegrees and then to degrees, gives $t(C_4,G)\ge t(K_2,G)^4$ for every graph $G$; random graphs come arbitrarily close to equality, but no graph with an edge attains it (\cref{sec:graphon:after-counting}). The optimizer is something that many sequences approach, and the finite theory has no name for it.

\subsection{Quasirandom graphs}

The random graph $G(n,p)$ satisfies $t(F,G(n,p))\to p^{e(F)}$ almost surely. \textcite{Thomason1987} and \textcite{ChungGrahamWilson1989} asked which deterministic properties make a graph look random, and found that many natural ones are equivalent.\sidenote{Chung, Graham, and Wilson coined the term \emph{quasi-random}; they worked at density $\tfrac12$, noting that every fixed density can be treated alike. Thomason called his random-like graphs \emph{jumbled}.} Let $e(S,T)$ be the number of pairs $(s,t)\in S\times T$ with $st\in E(G)$. For graphs $G_n$ on $n$ vertices with edge densities tending to $p$, each of these properties implies the others: the $4$-cycle density tends to $p^4$; every density $t(F,G_n)$ tends to $p^{e(F)}$; $e(S,T)=p\abs{S}\abs{T}+o(n^2)$ uniformly over vertex sets $S,T$; the top eigenvalue of the adjacency matrix is $pn+o(n)$ and the others are $o(n)$; most pairs of vertices have about $p^2n$ common neighbours. Such sequences are \term{quasirandom}.\index{quasirandom graph} No single graph is quasirandom, since an $o(1)$ means nothing for one graph, and the proof is a chain of implications, each a page of estimates with error terms, like this one.

\begin{aside}[A finite proof that $4$-cycles control discrepancy]
Let $G$ have $n$ vertices, adjacency matrix $A$, edge density $p=2e(G)/n^2$, degrees $d(v)$, and codegrees $c(u,w)=\abs{N(u)\cap N(w)}$, and suppose $t(C_4,G)=p^4+\eta$ with $0\le\eta\le1$. Three estimates do the work. First, $\graphonhom(C_4,G)=\sum_{u,w}c(u,w)^2$ and $\sum_{u,w}c(u,w)=\sum_v d(v)^2\ge p^2n^3$, so the Cauchy--Schwarz inequality gives $\sum_v d(v)^2\le n^3(p^2+\eta^{1/2})$, and hence $\sum_v(d(v)-pn)^2\le\eta^{1/2}n^3$. Second, expanding the square,
\[
  \sum_{u,w}\bigl(c(u,w)-p^2n\bigr)^2=n^4t(C_4,G)-2p^2n\sum_v d(v)^2+p^4n^4\le\eta n^4 .
\]
Third, for vertex sets $S,T$, the Cauchy--Schwarz inequality over $u\in S$ gives
\[
  \bigl(e(S,T)-p\abs{S}\abs{T}\bigr)^2\le n\sum_{v,v'\in T}\bigl(c(v,v')-p\,d(v)-p\,d(v')+p^2n\bigr),
\]
and the first two estimates bound the right side by $3\eta^{1/4}n^4$. So $\abs{e(S,T)-p\abs{S}\abs{T}}\le\sqrt3\,\eta^{1/8}n^2$ for all $S$ and $T$. Every inequality carries an error term, and every error term must be carried to the end.
\end{aside}

\subsection{Regularity}

The deepest tool of the finite theory says that every graph is approximately a bounded number of random-like bipartite graphs. Let $d(X,Y)=e(X,Y)/(\abs{X}\abs{Y})$. Disjoint vertex sets $X$ and $Y$ form an \term{$\eps$-regular pair} if $\abs{d(X',Y')-d(X,Y)}\le\eps$ whenever $X'\subseteq X$, $Y'\subseteq Y$, $\abs{X'}\ge\eps\abs{X}$, and $\abs{Y'}\ge\eps\abs{Y}$.

\begin{theorem}[Szemer\'edi's regularity lemma, \textcite{Szemeredi1978}]\label{thm:graphon:szemeredi}
For every $\eps>0$ and every integer $m\ge1$ there is an integer $M$ such that every graph with at least $m$ vertices has a partition of its vertex set into $k$ parts, with $m\le k\le M$ and part sizes differing by at most one, in which all but at most $\eps k^2$ of the pairs of parts are $\eps$-regular.
\end{theorem}

Szemer\'edi had used a bipartite form of the lemma in his 1975 proof that dense sets of integers contain long arithmetic progressions.\index{regularity lemma!Szemer\'edi's} The proof refines partitions to raise an energy, $\sum_{i,j}\abs{V_i}\abs{V_j}d(V_i,V_j)^2/n^2\in[0,1]$, by a fixed power of $\eps$ per round; since one round can cut a part into $2^{k-1}$ pieces, to be re-cut into equal parts, the number of parts grows as a tower of twos. \textcite{Gowers1997} showed that $M$ must indeed be a tower of height a power of $1/\eps$. Here is the basic counting estimate, with its proof, because its hypotheses will matter.

\begin{lemma}[Triangle counting]\label{lem:graphon:count-classical}
Let $X,Y,Z$ be disjoint vertex sets of a graph such that the three pairs among them are $\eps$-regular with densities $d_{XY},d_{XZ},d_{YZ}\ge2\eps$. Then the number of triangles with one vertex in each set is at least $(1-2\eps)(d_{XY}-\eps)(d_{XZ}-\eps)(d_{YZ}-\eps)\abs{X}\abs{Y}\abs{Z}$.
\end{lemma}

\begin{proof}
Let $X_Y$ be the set of vertices of $X$ with fewer than $(d_{XY}-\eps)\abs{Y}$ neighbours in $Y$. If $\abs{X_Y}\ge\eps\abs{X}$, then the pair $(X_Y,Y)$ would have density less than $d_{XY}-\eps$, contradicting the regularity of $(X,Y)$; so $\abs{X_Y}<\eps\abs{X}$. The same holds for $Z$, so at least $(1-2\eps)\abs{X}$ vertices $x\in X$ have at least $(d_{XY}-\eps)\abs{Y}$ neighbours in $Y$ and at least $(d_{XZ}-\eps)\abs{Z}$ neighbours in $Z$. These neighbourhoods have at least $\eps\abs{Y}$ and $\eps\abs{Z}$ elements, because $d_{XY},d_{XZ}\ge2\eps$; so the regularity of $(Y,Z)$ gives at least $(d_{YZ}-\eps)$ times the product of their sizes edges between them, and each such edge closes a triangle with $x$. Summing over these vertices $x$ gives the bound.
\end{proof}

The lemma counts from below only, and only when the densities are at least $2\eps$; for a general $F$ the idea needs an induction on $v(F)$ that keeps the sets above the threshold where regularity applies. Regularity and counting give the triangle removal lemma, implicit in the work of \textcite{RuzsaSzemeredi1978} on triple systems:\index{removal lemma} for every $\eps>0$ there is $\delta>0$ such that every graph on $n$ vertices with at most $\delta n^3$ triangles can be made triangle-free by deleting at most $\eps n^2$ edges. Applied to a graph built from a set $A\subseteq\{1,\dots,N\}$, it gives Roth's theorem of 1953: if $A$ has no three-term arithmetic progression, then $\abs{A}=o(N)$. The proof gives $1/\delta$ a tower of height polynomial in $1/\eps$.

Three features stand out. Its theorems are about sequences, and the objects they are about, \emph{the} random-like graph and \emph{the} optimal graph of density $p$, do not exist. Its lemmas carry hypotheses that only track error terms. And its numbers are towers. The completion removes the first two features and trades the third for something else (\cref{sec:graphon:costs}).

\section{The move: a graph is a function on the unit square}\label{sec:graphon:move}

Grothendieck's advice was to stoop low. Here the humble object is the adjacency matrix, drawn as a black-and-white picture on the unit square and allowed to blur into shades of gray. We write $\lambda$ for Lebesgue measure.

\begin{definition}[Graphons]\label{def:graphon:graphon}
A \term{kernel} is a bounded measurable function $U\colon[0,1]^2\to\R$, and a \term{graphon} is a kernel $W$ with values in $[0,1]$ that is symmetric: $W(x,y)=W(y,x)$. For a graph $F$ and a symmetric kernel $U$ the homomorphism density is
\[
  t(F,U)\defeq\int_{[0,1]^{V(F)}}\prod_{ij\in E(F)}U(x_i,x_j)\dd x .
\]
For a graph $G$ on $\{1,\dots,n\}$, split $[0,1]$ into the intervals $I_u={[(u-1)/n,\,u/n)}$, with the point $1$ added to $I_n$. The \term{step graphon} $W_G$ equals $1$ on $I_u\times I_v$ when $uv\in E(G)$, and $0$ elsewhere.
\end{definition}

\index{graphon}\index{step graphon}

\begin{lemma}\label{lem:graphon:step}
Every pair of graphs $F$, $G$ satisfies $t(F,W_G)=t(F,G)$.
\end{lemma}

\begin{proof}
Split $[0,1]^{V(F)}$ into $n^{v(F)}$ boxes, one for each map $\phi\colon V(F)\to V(G)$: the box of $\phi$ is the product of the intervals $I_{\phi(i)}$ over $i\in V(F)$. On this box the integrand equals $1$ if $\phi$ is a homomorphism and $0$ otherwise, and the box has volume $n^{-v(F)}$.
\end{proof}

\begin{figure}[t]
\centering
\begin{tikzpicture}[x=1mm,y=1mm]
  \def\sd{18}
  % top row: G_n with u ~ v iff u + v <= n, and the limit 1{x+y<=1}
  \foreach \u in {1,...,8} {\foreach \v in {1,...,8} {
    \pgfmathtruncatemacro{\s}{\u+\v}
    \ifnum\u=\v\relax\else\ifnum\s<9\relax
      \fill[sea] ({(\u-1)*\sd/8},{-(\v-1)*\sd/8}) rectangle ++({\sd/8},{-\sd/8});
    \fi\fi}}
  \foreach \u in {1,...,24} {\foreach \v in {1,...,24} {
    \pgfmathtruncatemacro{\s}{\u+\v}
    \ifnum\u=\v\relax\else\ifnum\s<25\relax
      \fill[sea] ({25+(\u-1)*\sd/24},{-(\v-1)*\sd/24}) rectangle ++({\sd/24},{-\sd/24});
    \fi\fi}}
  \fill[sea] (50,0) -- (50+\sd,0) -- (50,-\sd) -- cycle;
  % bottom row: Paley graphs P_13, P_29 and the constant 1/2
  \foreach \u in {0,...,12} {\foreach \r in {1,3,4,9,10,12} {
    \pgfmathtruncatemacro{\v}{mod(\u+\r,13)}
    \fill[sea] ({\u*\sd/13},{-26-\v*\sd/13}) rectangle ++({\sd/13},{-\sd/13});}}
  \foreach \u in {0,...,28} {\foreach \r in {1,4,5,6,7,9,13,16,20,22,23,24,25,28} {
    \pgfmathtruncatemacro{\v}{mod(\u+\r,29)}
    \fill[sea] ({25+\u*\sd/29},{-26-\v*\sd/29}) rectangle ++({\sd/29},{-\sd/29});}}
  \fill[sea!50] (50,-26) rectangle ++(\sd,-\sd);
  % frames, arrows, labels
  \foreach \oy in {0,-26} {\foreach \ox in {0,25,50} {
    \draw[muted,line width=0.3pt] (\ox,\oy) rectangle ++(\sd,-\sd);}
    \node[muted] at (46.5,\oy-9) {$\to$};}
  \node[below] at (9,-18) {\small $W_{G_8}$};
  \node[below] at (34,-18) {\small $W_{G_{24}}$};
  \node[below] at (59,-18) {\small $\ind_{\{x+y\le1\}}$};
  \node[below] at (9,-44) {\small $W_{P_{13}}$};
  \node[below] at (34,-44) {\small $W_{P_{29}}$};
  \node[below] at (59,-44) {\small $\tfrac12$};
\end{tikzpicture}
\caption{Graphs as pictures, with $x$ increasing to the right and $y$ downward, as in an adjacency matrix. Top: the graphs $G_n$ on $\{1,\dots,n\}$ in which $u\ne v$ are adjacent when $u+v\le n$, for $n=8$ and $24$, and their limit, the indicator of $\{x+y\le1\}$. Bottom: the Paley graphs $P_{13}$ and $P_{29}$, whose pictures converge in cut distance to the constant graphon $\tfrac12$, drawn in uniform half tone.}
\label{fig:graphon:pixels}
\end{figure}

\Cref{fig:graphon:pixels} shows two sequences. In the first, the pictures converge in $L^1$ to the indicator of a triangle, and with them all densities, since telescoping gives $\abs{t(F,U)-t(F,W)}\le e(F)\norm{U-W}_1$ for graphons. In the second, $P_q$ is the Paley graph on the residues modulo a prime $q\equiv1\pmod4$, two residues being adjacent when their difference is a nonzero square. Paley graphs are quasirandom (Exercise~\ref{ex:graphon:eig}), so $t(F,P_q)\to2^{-e(F)}$, the $F$-density of the constant function $\tfrac12$: the pictures turn gray. But the constant $\tfrac12$ is the picture of no graph. A limit has escaped from the space of graphs, as Dirichlet's function escapes from the Riemann integrable functions in \cref{ch:lebesgue}, and the move is the same: put the limits into the space.

Completion needs a metric, and the natural candidates fail.\sidenote{In \cref{ch:lebesgue} the $L^1$ completion was the right one; here it fails, since $L^1$ limits of $\{0,1\}$-valued functions stay $\{0,1\}$-valued.} Since $W_G$ takes only the values $0$ and $1$, every graph has $\norm{W_G-\tfrac12}_1=\tfrac12$, however its vertices are labeled: in $L^1$ no sequence of graphs approaches the gray square. At the other extreme, the graphons $W_n(x,y)=\tfrac12+\tfrac12\cos(2\pi n(x-y))$ converge weakly in $L^2$ to $\tfrac12$, yet $t(C_4,W_n)=\tfrac1{16}+\tfrac1{128}$ for every $n$ (expand $W_n$ in Fourier modes), while $t(C_4,\tfrac12)=\tfrac1{16}$: in the weak topology the $4$-cycle density is not continuous. The right norm lies in between, and the algorithmists had already found it.\sidenote{\textcite{FriezeKannan1999} introduced the cut norm of a matrix to design approximation algorithms for dense instances of problems such as MAX-CUT; see Exercise~\ref{ex:graphon:maxcut}.}

\begin{definition}[Cut norm and cut distance]\label{def:graphon:cut}
The \term{cut norm} of a kernel $U$ is
\[
  \graphoncut{U}\defeq\sup_{S,T}\Bigl\lvert\int_{S\times T}U(x,y)\dd x\dd y\Bigr\rvert,
\]
the supremum over measurable sets $S,T\subseteq[0,1]$. A \term{measure-preserving bijection} is a bijection $\phi\colon[0,1]\to[0,1]$ such that $\phi$ and $\phi^{-1}$ are measurable and preserve $\lambda$; for a graphon $W$ let $W^\phi(x,y)\defeq W(\phi(x),\phi(y))$. The \term{cut distance} of graphons $U$ and $W$ is $\graphondelta(U,W)\defeq\inf_\phi\graphoncut{U-W^\phi}$, the infimum over measure-preserving bijections. Graphons at cut distance $0$ are \term{weakly isomorphic}, and $\graphonspace$ is the set of graphons modulo weak isomorphism, with the metric $\graphondelta$. A graph $G$ is a point of $\graphonspace$ through $W_G$.
\end{definition}

\index{cut norm}\index{cut distance}\index{weak isomorphism}

Since measure-preserving bijections can be inverted and composed, $\graphondelta$ is symmetric and satisfies the triangle inequality. Relabeling the vertices of $G$ permutes the intervals $I_u$, so $\graphondelta$ compares graphs up to isomorphism, and graphs of different sizes too. For a kernel constant on the squares $I_u\times I_v$, the integral over $S\times T$ is bilinear in the fractions $n\lambda(S\cap I_u)$ and $n\lambda(T\cap I_v)$, so it is extremal when $S$ and $T$ are unions of intervals, that is, vertex sets. In particular,
\[
  \graphoncut{W_G-p}=\max_{S,T\subseteq V(G)}\frac{\abs{e(S,T)-p\abs{S}\abs{T}}}{n^2},
\]
and the discrepancy property of quasirandom graphs is exactly the statement $\graphoncut{W_{G_n}-p}\to0$. So the cut norm is weak enough that random graphs converge to constants; the next lemma is the first sign that it is strong enough for counting.

\begin{lemma}[Testing the cut norm with functions]\label{lem:graphon:functions}
Every kernel $U$ and all measurable $f,g\colon[0,1]\to[0,1]$ satisfy $\bigl\lvert\int\!\!\int U(x,y)f(x)g(y)\dd x\dd y\bigr\rvert\le\graphoncut{U}$.
\end{lemma}

\begin{proof}
Fix $g$ and let $h(x)=\int U(x,y)g(y)\dd y$. Since $0\le f\le1$, the integral $\int fh$ lies between $\int_{\{h<0\}}h$ and $\int_{\{h>0\}}h$. So the absolute value of $\int\!\!\int U(x,y)f(x)g(y)$ is at most that of $\int\!\!\int U(x,y)\ind_A(x)g(y)$, for $A=\{h<0\}$ or $A=\{h>0\}$. With $A$ fixed, the same argument in the variable $y$ replaces $g$ by the indicator of a set $B$, and $\abs{\int_{A\times B}U}\le\graphoncut{U}$.
\end{proof}

The theory rests on four facts about the space $\graphonspace$.\index{compactness of graphon space}

\begin{theorem}[Graph limits]\label{thm:graphon:main}
\leavevmode
\begin{enumerate}
\item (Counting.) Every graph $F$ and all graphons $U$, $W$ satisfy $\abs{t(F,U)-t(F,W)}\le e(F)\,\graphondelta(U,W)$.
\item (Compactness.) The metric space $(\graphonspace,\graphondelta)$ is compact.
\item (Uniqueness.) If $t(F,U)=t(F,W)$ for every graph $F$, then $\graphondelta(U,W)=0$.
\item (Density.) Every graphon is the $\graphondelta$-limit of the step graphons of a sequence of graphs.
\end{enumerate}
\end{theorem}

Part (i), and the convergence of random graphs to their graphon, which gives (iv), are due to \textcite{LovaszSzegedy2006}, (ii) to \textcite{LovaszSzegedy2007}, and (iii) to the theorem of \textcite{BorgsEtAl2008} that convergence of all densities is equivalent to convergence in cut distance. By (ii) and (iv), $\graphonspace$ is the completion of the graphs under $\graphondelta$. By (i)--(iii), $W\mapsto(t(F,W))_F$ is a continuous injection of a compact space, hence a homeomorphism onto its image: convergence in cut distance is convergence of every density, and graphs whose densities all converge have a limit graphon, unique up to weak isomorphism. We prove (i) and sketch (ii) below; (iii) rests on a sampling argument that we omit \parencite[see][]{BorgsEtAl2008,Lovasz2012}, and (iv) is Exercise~\ref{ex:graphon:dense}.

In hindsight, three observations make the move. Densities are integrals of products of the picture $W_G$, so the candidates for limits are functions. A diagonal argument over the countably many graphs $F$ shows that every sequence of graphs has a subsequence along which all densities converge, so the question was never whether limits exist but what they are. And the norm had to fit the counting: weak enough that random graphs converge, strong enough that densities are continuous, and blind to relabeling. Lov\'asz and Szegedy posted the theory in 2004;\sidenote{Their construction of the limit used weak regular partitions and the martingale convergence theorem, the two ingredients of the proof of compactness sketched in \cref{sec:graphon:after-compact}.} \textcite{BorgsEtAl2008} related the cut distance to sampling and property testing, and \textcite{LovaszSzegedy2007} recast regularity lemmas as compactness in Hilbert space.

\section{After: counting, rigidity, and quasirandomness}\label{sec:graphon:after-counting}

The first dividend is that counting needs no hypotheses. To cover regular pairs as well as graphons, the counting lemma allows a different kernel on each edge of $F$.\index{counting lemma}

\begin{lemma}[Counting lemma]\label{lem:graphon:counting}
Let $F$ be a graph whose edges are listed as ordered pairs $(a_1,b_1),\dots,(a_m,b_m)$, and let $U_1,\dots,U_m$ and $W_1,\dots,W_m$ be kernels with values in $[0,1]$. Write $t(F;U_1,\dots,U_m)\defeq\int\prod_{j=1}^mU_j(x_{a_j},x_{b_j})\dd x$. Then
\[
  \bigl\lvert t(F;U_1,\dots,U_m)-t(F;W_1,\dots,W_m)\bigr\rvert\le\sum_{j=1}^m\graphoncut{U_j-W_j}.
\]
In particular, every pair of graphons $U$, $W$ satisfies $\abs{t(F,U)-t(F,W)}\le e(F)\,\graphondelta(U,W)$.
\end{lemma}

\begin{proof}
Replace the kernels one at a time. For $0\le j\le m$ let $\Theta_j$ be the density computed with $W_i$ on the first $j$ edges and $U_i$ on the others. The difference of the two densities is $\sum_{j=1}^m(\Theta_{j-1}-\Theta_j)$, and $\Theta_{j-1}-\Theta_j$ is the integral of $(U_j-W_j)(x_{a_j},x_{b_j})$ times $m-1$ factors with values in $[0,1]$. Fix all coordinates other than $x_{a_j}$ and $x_{b_j}$. Because $F$ has no multiple edges, each remaining factor depends on at most one of $x_{a_j}$ and $x_{b_j}$, so their product has the form $f(x_{a_j})g(x_{b_j})c$ with $f$, $g$, $c$ taking values in $[0,1]$. By \cref{lem:graphon:functions} the integral over $x_{a_j}$ and $x_{b_j}$ is at most $\graphoncut{U_j-W_j}$ in absolute value, and integrating over the other coordinates keeps this bound.

For graphons, let $\phi$ be a measure-preserving bijection. The change of variables $x_i\mapsto\phi(x_i)$ in every coordinate shows $t(F,W^\phi)=t(F,W)$, so the first part, with $U_j=U$ and $W_j=W^\phi$, gives $\abs{t(F,U)-t(F,W)}\le e(F)\graphoncut{U-W^\phi}$. Take the infimum over $\phi$.
\end{proof}

The classical counting lemma follows by encoding vertex sets as kernels; the only work is the extreme-point remark made after \cref{def:graphon:cut}, applied to one pair of sets.

\begin{corollary}[Counting in regular pairs]\label{cor:graphon:count-regular}
Let $F$ be a graph on $\{1,\dots,k\}$, let $G$ be a graph, and let $V_1,\dots,V_k$ be disjoint nonempty subsets of $V(G)$ such that $(V_i,V_j)$ is $\eps$-regular for every edge $ij$ of $F$. Then the number $N$ of tuples $(v_1,\dots,v_k)\in V_1\times\dots\times V_k$ with $v_iv_j\in E(G)$ for every edge $ij$ of $F$ satisfies
\[
  \Bigl\lvert\frac{N}{\abs{V_1}\cdots\abs{V_k}}-\prod_{ij\in E(F)}d(V_i,V_j)\Bigr\rvert\le e(F)\,\eps .
\]
\end{corollary}

\begin{proof}
For each $i$, give each vertex of $V_i$ an interval of length $1/\abs{V_i}$, so that these intervals partition $[0,1]$. For an edge $ij$ of $F$ with $i<j$, let $A_{ij}(s,t)=1$ if the vertices of $V_i$ and $V_j$ whose intervals contain $s$ and $t$ are adjacent, and $A_{ij}(s,t)=0$ otherwise. Then $N/(\abs{V_1}\cdots\abs{V_k})=t(F;(A_{ij}))$, and the product of densities is $t(F;(d(V_i,V_j)))$ with constant kernels, so by \cref{lem:graphon:counting} it suffices to show $\graphoncut{A_{ij}-d}\le\eps$ for $d=d(V_i,V_j)$. The kernel $A_{ij}-d$ is constant on products of intervals, so we may take $S$ and $T$ to be vertex sets $X\subseteq V_i$ and $Y\subseteq V_j$, for which $\int_{S\times T}(A_{ij}-d)=(e(X,Y)-d\abs{X}\abs{Y})/(\abs{V_i}\abs{V_j})$. If $\abs{X}\ge\eps\abs{V_i}$ and $\abs{Y}\ge\eps\abs{V_j}$, then regularity bounds the absolute value of the numerator by $\eps\abs{X}\abs{Y}\le\eps\abs{V_i}\abs{V_j}$. Otherwise, say $\abs{X}<\eps\abs{V_i}$; then $e(X,Y)$ and $d\abs{X}\abs{Y}$ both lie in $[0,\abs{X}\abs{Y}]$, so their difference is less than $\eps\abs{V_i}\abs{V_j}$.
\end{proof}

For triangles, the corollary counts $(d_{XY}d_{XZ}d_{YZ}\pm3\eps)\abs{X}\abs{Y}\abs{Z}$ triangles in a regular triple: an upper bound as well as a lower one, with no lower bound on the densities, and with one proof for every $F$ instead of an induction. The price is an additive error, useless when the product of densities is below $e(F)\eps$; in the removal lemma this only changes how small $\eps$ must be.

\subsection{Goodman's bound and the rigidity of the constant}

\begin{theorem}[Goodman's bound for graphons]\label{thm:graphon:goodman}
Every graphon $W$ satisfies $t(K_3,W)\ge t(K_2,W)\bigl(2t(K_2,W)-1\bigr)$.
\end{theorem}

\begin{proof}
Let $p=t(K_2,W)$ and let $d(x)=\int W(x,y)\dd y$ be the degree function. Because $W$ takes values in $[0,1]$, we have $(1-W(x,z))(1-W(y,z))\ge0$, that is, $W(x,z)W(y,z)\ge W(x,z)+W(y,z)-1$. Integrate over $z$, multiply by $W(x,y)\ge0$, and integrate over $x$ and $y$; by the symmetry of $W$, both $\int\!\!\int W(x,y)d(x)$ and $\int\!\!\int W(x,y)d(y)$ equal $\int d^2$, so
\[
  t(K_3,W)\ge2\int d^2-p\ge2p^2-p .
\]
The last step is the Cauchy--Schwarz inequality $\int d^2\ge(\int d)^2=p^2$.
\end{proof}

Applied to $W_G$, this is \cref{thm:graphon:goodman-finite}, by the same steps without the bookkeeping. The proof is not much shorter, and that is not the point: the inequality now lives on a compact space, where one can ask whether it is attained, and by what. For the $4$-cycle the answer is sharp.

\begin{lemma}[Rigidity of the constant graphon]\label{lem:graphon:rigidity}
Every graphon $W$ satisfies $t(C_4,W)\ge t(K_2,W)^4$, with equality only if $W$ is almost everywhere constant.
\end{lemma}

\begin{proof}
The plan is to apply the Cauchy--Schwarz inequality twice, as in the finite proof, and then to read off what equality forces. Let $p=t(K_2,W)$, let $d$ be the degree function, and let $c(x,y)=\int W(x,z)W(y,z)\dd z$ be the codegree function. For the $4$-cycle with vertices $x_1,x_2,x_3,x_4$ in cyclic order, integrating over $x_2$ and $x_4$ first gives $t(C_4,W)=\int\!\!\int c(x_1,x_3)^2$, and integrating over $x$ and $y$ first gives $\int\!\!\int c=\int d^2$. On a probability space $\int g^2\ge(\int g)^2$, with equality only when $g$ is almost everywhere constant. Apply this to $c$ and then to $d$:
\[
  t(C_4,W)=\int\!\!\int c^2\ge\Bigl(\int\!\!\int c\Bigr)^2=\Bigl(\int d^2\Bigr)^2\ge\Bigl(\int d\Bigr)^4=p^4 .
\]
Suppose now that equality holds. Then $c=p^2$ and $d=p$ almost everywhere. Let $T_W$ be the integral operator $T_Wf(x)=\int W(x,y)f(y)\dd y$ on $L^2[0,1]$. For $f$ with $\int f=0$, Fubini's theorem gives $\int(T_Wf)^2=\int\!\!\int f(y)f(y')\,c(y,y')\dd y\dd y'=p^2(\int f)^2=0$, so $T_Wf=0$. For measurable sets $A,B\subseteq[0,1]$, write $\ind_B=\lambda(B)+f$ with $\int f=0$; then $T_W\ind_B=\lambda(B)\,d=p\lambda(B)$, and integrating over $A$ gives $\int_{A\times B}W=p\lambda(A)\lambda(B)$. So the finite measures with densities $W$ and $p$ on $[0,1]^2$ agree on measurable rectangles, a class closed under intersection that generates the Borel sets; by the $\pi$--$\lambda$ theorem they agree on all Borel sets, and $W=p$ almost everywhere.
\end{proof}

The inequality is the finite one of \cref{sec:graphon:before}; what is new is the equality case. Since $W_G$ vanishes on the squares $I_u\times I_u$, it is not almost everywhere constant unless $G$ has no edges, so no graph with an edge attains equality. The equality case is exactly what quasirandomness needs.

\begin{theorem}[Quasirandom graphs, after \textcite{ChungGrahamWilson1989}]\label{thm:graphon:cgw}
Let $p\in[0,1]$, and let $(G_n)$ be a sequence of graphs. The following are equivalent:
\begin{enumerate}
\item $t(K_2,G_n)\to p$ and $t(C_4,G_n)\to p^4$;
\item $t(F,G_n)\to p^{e(F)}$ for every graph $F$;
\item $\graphoncut{W_{G_n}-p}\to0$, that is, $G_n$ converges to the constant graphon $p$.
\end{enumerate}
\end{theorem}

\begin{proof}
Statement (iii) implies (ii) by the counting lemma, since the constant graphon $p$ has $F$-density $p^{e(F)}$, and (ii) implies (i) by taking $F=K_2$ and $F=C_4$. Suppose (i) holds and (iii) fails. Then for some $\eps>0$ a subsequence has $\graphoncut{W_{G_n}-p}\ge\eps$. By compactness (\cref{thm:graphon:main}(ii)), a further subsequence converges in cut distance to a graphon $W$. By the counting lemma, $t(K_2,W)=p$ and $t(C_4,W)=p^4$, so $W=p$ almost everywhere by \cref{lem:graphon:rigidity}. The constant graphon is fixed by every measure-preserving bijection, so $\graphondelta(W_{G_n},p)=\graphoncut{W_{G_n}-p}$, and this tends to $0$ along the further subsequence, a contradiction.
\end{proof}

Statement (iii) is the discrepancy property, by the formula after \cref{def:graphon:cut}; the eigenvalue property joins the list in the same way (Exercise~\ref{ex:graphon:eig}), as does the codegree property, and no growth of $v(G_n)$ is assumed. The chain of implications has become one rigidity lemma and one appeal to compactness. What the new proof does not give is a rate (\cref{sec:graphon:costs}).

\section{After: regularity, compactness, and extremal graphons}\label{sec:graphon:after-compact}

The regularity lemma reappears in the completion as a statement about projections in Hilbert space, with single-exponential bounds.

\begin{definition}[Stepping]\label{def:graphon:stepping}
Let $P=\{S_1,\dots,S_k\}$ be a partition of $[0,1]$ into measurable sets. For a kernel $W$, the \term{stepping} $W_P$ is the kernel whose value on $S_i\times S_j$ is the average $\lambda(S_i)^{-1}\lambda(S_j)^{-1}\int_{S_i\times S_j}W$; parts of measure zero play no role, and we ignore them. Let $L_P$ be the space of kernels that are constant on each product $S_i\times S_j$.
\end{definition}

The stepping $W_P$ is the orthogonal projection of $W$ onto $L_P$ in $L^2([0,1]^2)$, since it lies in $L_P$ and $W-W_P$ integrates to zero over each $S_i\times S_j$; in probabilistic language, it is a conditional expectation. If $P'$ refines $P$, then $L_P\subseteq L_{P'}$ and $(W_{P'})_P=W_P$.

\begin{lemma}[Weak regularity, \textcite{FriezeKannan1999}]\label{lem:graphon:weakreg}
Let $W$ be a graphon, let $\eps>0$, and let $P_0$ be a partition of $[0,1]$ into finitely many measurable sets. There is a partition $P$ refining $P_0$, with at most $4^{1/\eps^2}\abs{P_0}$ parts, such that $\graphoncut{W-W_P}\le\eps$.
\end{lemma}

\index{regularity lemma!weak}

\begin{proof}
The plan is to refine $P_0$ step by step, each step raising the \emph{energy} $\norm{W_P}_2^2$ by more than $\eps^2$; the energy never exceeds $\norm{W}_2^2\le1$, because projections do not increase norms.

Suppose that a partition $P$ has $\graphoncut{W-W_P}>\eps$, and choose measurable $S$ and $T$ with $\abs{\int_{S\times T}(W-W_P)}>\eps$. Let $P'$ be the common refinement of $P$, $\{S,[0,1]\setminus S\}$, and $\{T,[0,1]\setminus T\}$; each part of $P$ splits into at most four parts of $P'$. The indicator $\ind_{S\times T}$ lies in $L_{P'}$, so it is orthogonal to $W-W_{P'}$, and therefore $\ip{W_{P'}-W_P}{\ind_{S\times T}}=\ip{W-W_P}{\ind_{S\times T}}$, whose absolute value exceeds $\eps$. By the Cauchy--Schwarz inequality and $\norm{\ind_{S\times T}}_2\le1$, this gives $\norm{W_{P'}-W_P}_2>\eps$. Since $W_P$ is the projection of $W_{P'}$ onto $L_P$, Pythagoras' theorem gives $\norm{W_{P'}}_2^2=\norm{W_P}_2^2+\norm{W_{P'}-W_P}_2^2>\norm{W_P}_2^2+\eps^2$.

Repeat this step, starting from $P_0$, while the cut norm exceeds $\eps$. After $s$ steps the energy exceeds $s\eps^2$, so fewer than $1/\eps^2$ steps occur, and the final partition has at most $4^{1/\eps^2}\abs{P_0}$ parts.
\end{proof}

This is Szemer\'edi's energy increment without irregular pairs to track or parts to equalize, so the number of parts is single exponential, and \textcite{ConlonFox2012} showed that $2^{\Omega(\eps^{-2})}$ parts are sometimes necessary. The conclusion is weaker, but it is all that counting and compactness need; for graphs it gives Frieze and Kannan's partition (Exercise~\ref{ex:graphon:maxcut}).

\begin{theorem}[Compactness, \textcite{LovaszSzegedy2007}]\label{thm:graphon:compact}
Every sequence of graphons has a subsequence that converges in cut distance to a graphon. Equivalently, $(\graphonspace,\graphondelta)$ is compact.
\end{theorem}

\begin{proofsketch}
The plan is to approximate each $W_n$ by step graphons of bounded complexity at every scale, take limits scale by scale, and glue the scales with the martingale convergence theorem. Two measure-theoretic points are omitted; see \textcite{LovaszSzegedy2007} and \textcite[Chapter~4]{Zhao2023}.

\emph{Regularize.} Applying \cref{lem:graphon:weakreg} with $\eps=1/k$ to the previous partition gives, for each $n$, nested partitions $P_{n,1},P_{n,2},\dots$ such that $W_{n,k}\defeq(W_n)_{P_{n,k}}$ satisfies $\graphoncut{W_n-W_{n,k}}\le1/k$; allowing empty parts, $P_{n,k}$ has exactly $m_k\defeq\prod_{j\le k}4^{j^2}$ parts.

\emph{Straighten.} Laying out the parts of $P_{n,1}$ as consecutive intervals, those of $P_{n,2}$ as consecutive subintervals of these, and so on, gives nested interval partitions $J_{n,k}$ and step graphons $U_{n,k}$ with steps $J_{n,k}$ and the values of $W_{n,k}$. Then $U_{n,k}=(U_{n,k+1})_{J_{n,k}}$, and $\graphondelta(W_{n,k},U_{n,k})=0$ by a measure-preserving bijection that we do not construct (the first omitted point).

\emph{Diagonalize.} Each $U_{n,k}$ is described by the nesting of its steps, which takes finitely many forms, and $m_k+m_k^2$ numbers in $[0,1]$, so a diagonal argument gives a subsequence along which, for every $k$, $U_{n,k}$ converges almost everywhere to a step graphon $U_k$; its interval steps $J_k$ are nested in $k$, and $U_k=(U_{k+1})_{J_k}$.

\emph{Glue.} On the probability space $[0,1]^2$ the $U_k$ form a martingale for the $\sigma$-algebras generated by the products of intervals of $J_k$ (intervals of length zero need care, the second omitted point). It takes values in $[0,1]$, so it converges almost everywhere and in $L^1$ to a graphon $U$.

\emph{Conclude.} For every $k$ and every $n$ in the subsequence,
\[
  \graphondelta(W_n,U)\le\graphoncut{W_n-W_{n,k}}+\graphondelta(W_{n,k},U_{n,k})+\norm{U_{n,k}-U_k}_1+\norm{U_k-U}_1 .
\]
The first term is at most $1/k$, the second is $0$, the third tends to $0$ as $n\to\infty$ by dominated convergence, and the fourth tends to $0$ as $k\to\infty$.
\end{proofsketch}

\subsection{New reach: extremal graphons}

\begin{corollary}[Extremal graphons exist]\label{cor:graphon:extremal}
Let $p\in[0,1]$ and let $F$ be a graph. The minimum of $t(F,W)$ over graphons $W$ with $t(K_2,W)=p$ is attained, and it is the least value of $\liminf_n t(F,G_n)$ over sequences of graphs with $t(K_2,G_n)\to p$.
\end{corollary}

\begin{proof}
By the counting lemma, $t(K_2,\cdot)$ and $t(F,\cdot)$ are continuous on $\graphonspace$, so the set where $t(K_2,\cdot)=p$ is compact (\cref{thm:graphon:compact}), and it contains the constant $p$; hence $t(F,\cdot)$ attains a minimum on it. If $t(K_2,G_n)\to p$, then every subsequence of $(G_n)$ has a further subsequence converging to a point of this set, so $\liminf_nt(F,G_n)$ is at least the minimum; and part~(iv) of \cref{thm:graphon:main} gives graphs converging to a minimizer.
\end{proof}

For the $4$-cycle the minimizer is the constant graphon, and by \cref{lem:graphon:rigidity} it is the only one: the optimizer that did not exist among graphs exists, and is unique, in the completion. For the triangle the minimum is the lower boundary of the region in \cref{fig:graphon:region}, a well-known open problem until Razborov solved it.\index{Razborov's triangle density theorem}

\begin{theorem}[\textcite{Razborov2008}]\label{thm:graphon:razborov}
Let $t\ge1$ be an integer, let $p\in[1-1/t,\,1-1/(t+1)]$, and let $s=\sqrt{t(t-(t+1)p)}$. The minimum of $t(K_3,W)$ over graphons $W$ with $t(K_2,W)=p$ is
\[
  \frac{(t-1)(t-2s)(t+s)^2}{t^2(t+1)^2},
\]
attained by the complete $(t+1)$-partite graphon with $t$ parts of measure $c=(t+s)/(t(t+1))$ and one part of measure $1-tc$.
\end{theorem}

Here a complete multipartite graphon with parts of measures $a_i$ is $1$ on $I_i\times I_j$ for $i\ne j$ and $0$ on $I_i\times I_i$, for consecutive intervals $I_i$ of lengths $a_i$. For $t=1$ the formula gives $0$, and at $p=1-1/r$ it agrees with Goodman's bound. Strictly inside each interval the minimizers are not unique: the complete bipartite graphon between the two smallest parts can be replaced by every triangle-free graphon on their union with the same integral, since each of its edges lies in the same number of triangles.

\begin{figure}[t]
\centering
\begin{tikzpicture}[x=62mm,y=34mm,
  declare function={rs(\p,\t)=sqrt(max(0,\t*(\t-\p*(\t+1))));
    raz(\p,\t)=((\t-1)/\t)*((\t-2*rs(\p,\t))/\t)*((\t+rs(\p,\t))/(\t+1))^2;}]
  \fill[seafoam] (0,0) -- plot[domain=0:1,samples=80] (\x,{\x*sqrt(\x)}) -- (1,0) -- cycle;
  \fill[white] (0.5,0) foreach \t in {2,...,14} {
      -- plot[domain={1-1/\t}:{1-1/(\t+1)},samples=14] (\x,{raz(\x,\t)}) } -- (1,1) -- (1,0) -- cycle;
  \draw[muted,line width=0.3pt] (0,0) rectangle (1,1);
  \draw[muted,dotted,line width=0.7pt] plot[domain=0:1,samples=50] (\x,{\x*\x*\x});
  \draw[seadeep,dashed,line width=0.5pt] plot[domain=0.5:1,samples=50] (\x,{\x*(2*\x-1)});
  \draw[sea,line width=0.9pt] plot[domain=0:1,samples=80] (\x,{\x*sqrt(\x)});
  \draw[seadeep,line width=0.9pt] (0,0) -- (0.5,0) foreach \t in {2,...,14} {
      -- plot[domain={1-1/\t}:{1-1/(\t+1)},samples=14] (\x,{raz(\x,\t)}) } -- (1,1);
  \foreach \r in {2,...,7} {\fill[seadeep] ({1-1/\r},{(1-1/\r)*(1-2/\r)}) circle[radius=0.7pt];}
  \node[below] at (0,0) {\small $0$};
  \node[below] at (0.5,0) {\small $\tfrac12$};
  \node[below] at (1,0) {\small $1$};
  \node[left] at (0,0) {\small $0$};
  \node[left] at (0,1) {\small $1$};
  \node[below] at (0.75,-0.08) {\small $t(K_2,W)$};
  \node[left] at (0,0.62) {\small $t(K_3,W)$};
  \node[sea] at (0.36,0.33) {\small $p^{3/2}$};
  \node[muted] at (0.37,0.12) {\small $p^{3}$};
  \draw[seadeep,line width=0.3pt] (0.80,0.27) -- (0.775,0.42);
  \node[seadeep,below] at (0.80,0.28) {\small $p(2p-1)$};
\end{tikzpicture}
\caption{The edge--triangle region: the pairs $(t(K_2,W),t(K_3,W))$ for graphons $W$ fill the shaded region (Exercise~\ref{ex:graphon:region}). The upper boundary $p^{3/2}$ is attained by cliques. The lower boundary is Razborov's scalloped curve (\cref{thm:graphon:razborov}), which meets Goodman's parabola $p(2p-1)$ (dashed) at the points $p=1-1/r$ (dots). The dotted curve $p^3$ belongs to the constant graphons, the limits of random and quasirandom graphs.}
\label{fig:graphon:region}
\end{figure}

How did Razborov prove the lower bound? \textcite{Razborov2007} built a calculus for exactly the objects of this chapter.\index{flag algebra} Its \emph{flag algebras} consist of formal combinations of densities of small graphs with some vertices labeled, with a product and an averaging operation compatible with limits. An inequality is proved by a sum of squares of such combinations, found by a semidefinite program, whose average equals the difference of the two sides up to visibly nonnegative terms; the certificate holds on the compact space of graphons, hence for all large graphs. Its first application reproved a partial result of Fisher on the triangle problem,\sidenote{Flag algebras have since become a standard tool of extremal combinatorics. A flag-algebra proof is a sum-of-squares proof carried out in the completion, the paradigm of \cref{ch:sos}.} and \cref{thm:graphon:razborov} followed.

\section{Costs and limits}\label{sec:graphon:costs}

\paragraph{Where did the difficulty go?} Into the compactness theorem, and through it into the weak regularity lemma, the martingale convergence theorem, and the measure theory of relabelings; the counting lemma is cheap. The cost is paid once, and each application costs a few lines. The catch is that compactness is ineffective. The proof of \cref{thm:graphon:cgw} says nothing about how small $\graphoncut{W_G-p}$ is when $t(C_4,G)=p^4+\eta$, while the finite proof in the aside gives $\sqrt3\,\eta^{1/8}$; a bound in the limit language needs the finite work again (Exercise~\ref{ex:graphon:quantitative}). Limit proofs of removal lemmas give no bounds; the improvement of the triangle removal lemma from a tower of height polynomial in $1/\eps$ to one of height $O(\log(1/\eps))$ came from combinatorics \parencite{Fox2011}. Nor did extremal problems become easy: the completion made Razborov's problem well posed, but solving it took flag algebras.

\paragraph{What did not dissolve.} Sparse graphs, first: if $e(G_n)=o(v(G_n)^2)$, then $\graphoncut{W_{G_n}}\le\norm{W_{G_n}}_1\to0$, so every sparse sequence converges to the zero graphon. Graphs of bounded degree need a different completion, by local limits \parencite[see][]{Lovasz2012}. For sparse graphs with growing degrees, the unbounded graphons of \textcite{BorgsChayesCohnZhao2019} keep compactness, under a hypothesis limiting dense spots, but lose counting: deleting the edges in triangles of $G(n,\rho_n)$, with $n\rho_n\to\infty$ and $\rho_n\sqrt n\to0$, moves the normalized cut distance by $o(1)$ and the normalized triangle density from $1$ to $0$. Second, nothing solves all extremal problems: \textcite{HatamiNorine2011} proved that it is undecidable whether a linear inequality between homomorphism densities holds for all graphs, and that some valid inequalities have no sum-of-squares certificate. Third, Sidorenko's conjecture \parencite{Sidorenko1993}\index{Sidorenko's conjecture}, that $t(H,W)\ge t(K_2,W)^{e(H)}$ for every bipartite $H$ and every graphon $W$, has the cleanest possible integral form and resisted proof for decades, though it holds for trees, even cycles, hypercubes, and many other families. A preprint of September 2026 \parencite{OpenAI2026} claims a counterexample with $35$ vertices and $66$ edges, built as a kernel with a strict density gap from which a finite graph is then sampled, the move of this chapter run in reverse. When this chapter was written the claim had not been refereed or independently checked. Fourth, a function of three variables does not determine the densities of $3$-uniform hypergraphs (Exercise~\ref{ex:graphon:hypergraphs}); \textcite{ElekSzegedy2012} needed variables for pairs of vertices as well.

\paragraph{Other seas.} Probability offers the same completion. The \term{$W$-random graph} $G(n,W)$ joins $i$ and $j$, independently, with probability $W(x_i,x_j)$, for independent uniform labels $x_i$, and $G(n,W)\to W$ with probability $1$ \parencite{LovaszSzegedy2006}. On the vertex set $\N$ it is an \emph{exchangeable} random graph, and around 1980 the probabilist Aldous and the logician Hoover proved that every exchangeable random graph is a mixture of $W$-random graphs, a two-dimensional de Finetti theorem; \textcite{DiaconisJanson2008} made the dictionary with graph limits explicit.\index{exchangeable random graph} The completion had been in probability theory all along, as a representation theorem, and \cref{ch:conformal} exploits the one-dimensional case of the same symmetry. Ultraproducts give limits without choosing a norm, which is how \textcite{ElekSzegedy2012} reached hypergraphs. Model theory gives a different regularity, for graphs without large half-graphs (\cref{ch:littlestone}). And the move completes other structures: words, permutations, and graphs carrying data (Exercises~\ref{ex:graphon:words}--\ref{ex:graphon:permutons} and~\ref{ex:graphon:gnn}).

\section{The move, distilled}\label{sec:graphon:distilled}

\begin{move}{Completion}
\movepart{In one sentence} When every theorem about large objects is a statement about sequences, find the norm in which the quantities you count are continuous and the space is compact, and let the sequences become points.
\movepart{Triggers} Theorems that hold ``for $n$ sufficiently large'' or ``up to $o(1)$''; extremal problems whose optimizers are sequences; a lemma saying that every object is close to one of boundedly many templates, a compactness theorem in disguise.
\movepart{Recipe} Encode each finite object as a function on a fixed probability space, and the quantities of interest as integrals. Choose a norm weak enough that natural sequences converge and strong enough that the quantities are continuous, and divide out the symmetries of the encoding. Prove compactness and the density of the finite objects once. Then prove theorems in the limit, where $o(1)$ becomes $0$ and inequalities acquire equality cases, and transfer them back by compactness.
\movepart{Anti-pattern} Completing in the wrong norm: in $L^1$ random graphs converge to nothing, and in the weak topology the $4$-cycle density is not continuous. Or a ``limit object'' without a compactness theorem: rewriting finite arguments in integral notation, with every $\eps$ still in place, is a renaming.
\end{move}

\dissolutiontest{Counting becomes one telescoping sum instead of an induction; quasirandomness becomes the rigidity of the constant graphon plus compactness, instead of chains of $o(1)$ estimates.}{Counting needs no lower bound on pair densities and covers every $F$; quasirandomness needs no growth of $v(G_n)$; extremal problems need no ``for $n$ sufficiently large''.}{The Chung--Graham--Wilson properties become one statement, convergence to a constant; regularity, sampling, and limits become facets of one compact space; graph limits and exchangeable random graphs become the same objects.}{Extremal graphons exist, and extremal problems become optimization over a compact space, the setting of flag algebras and Razborov's theorem. Rates and bounds are not part of the reach.}

This chapter and \cref{ch:lebesgue} make the same move with a difference: Lebesgue completed in a norm given in advance, while here the norm itself was the discovery, chosen to make counting continuous, and the quotient by relabeling was as necessary as the norm. Compactness, which the integral did not need, is supplied for permutations by the spaces of measures of \cref{ch:transport} (Exercise~\ref{ex:graphon:permutons}). Regularity returns in \cref{ch:littlestone}: graphs without long half-graphs, those of bounded Littlestone dimension, have regular partitions with no irregular pairs. Flag algebras are the sum-of-squares paradigm of \cref{ch:sos} applied in the completion, and the step from eigenvalues to discrepancy in Exercise~\ref{ex:graphon:eig} is the expander mixing lemma of \cref{ch:pseudorandom}.

\begin{exercises}
\item \consolidation\label{ex:graphon:region} \emph{The edge--triangle region.} (a) The step graphon of $K_r$ is the complete $r$-partite graphon with equal parts; show that its edge and triangle densities are $1-1/r$ and $(1-1/r)(1-2/r)$, so that Goodman's bound is attained. (b) Show that the clique graphon $\ind_{[0,a]^2}$ has $t(K_3)=t(K_2)^{3/2}$, and that $t(K_3,W)\le t(K_2,W)^{3/2}$ for every graphon $W$. (c) Show that for each $p$ the values $t(K_3,W)$ with $t(K_2,W)=p$ form a closed interval, which gives, with \cref{thm:graphon:razborov}, the region of \cref{fig:graphon:region}. \hint{For (b), write $t(K_3,W)=\int\!\!\int Wc$ with the codegree function $c$, and use the Cauchy--Schwarz inequality, $c(x,y)^2\le c(x,x)\,c(y,y)$, and $W^2\le W$. For (c), interpolate linearly.}

\item \consolidation\hard\label{ex:graphon:quantitative} \emph{Getting a bound back.} (a) Show that every symmetric kernel $U$ with values in $[-1,1]$ has $\graphoncut{U}\le t(C_4,U)^{1/4}$. (b) Let $W$ be a graphon whose degree function equals $p$ almost everywhere. Show that $t(C_4,W-p)=t(C_4,W)-p^4$, so that $\graphoncut{W-p}\le(t(C_4,W)-p^4)^{1/4}$. Why does the proof of \cref{thm:graphon:cgw} give no such bound? \hint{For (a), apply the Cauchy--Schwarz inequality in $x$, and then in the pair $(y,y')$ against the codegree function of $U$.}

\item \consolidation\label{ex:graphon:eig} \emph{Eigenvalues.} Let $G_n$ have $n$ vertices, adjacency eigenvalues $\lambda_1\ge\dots\ge\lambda_n$, and $t(K_2,G_n)\to p$. (a) Show that $t(C_4,G_n)=\sum_i\lambda_i^4/n^4$ and $\lambda_1\ge t(K_2,G_n)\,n$, and deduce that $t(C_4,G_n)\to p^4$ if and only if $\lambda_1/n\to p$ and $\max_{i\ge2}\abs{\lambda_i}/n\to0$. (b) The Paley graph $P_q$ has eigenvalues $(q-1)/2$ and $(-1\pm\sqrt q)/2$; conclude that $P_q$ converges to the constant $\tfrac12$. (c) Show that the eigenvalue condition alone, without $t(K_2,G_n)\to p$, does not imply quasirandomness.

\item \consolidation\label{ex:graphon:maxcut} \emph{Weak regularity for graphs, and MAX-CUT.} (a) Let $G$ be a graph on $\{1,\dots,n\}$. Show that if $P$ partitions $[0,1]$ into unions of intervals $I_u$, the supremum defining $\graphoncut{W_G-(W_G)_P}$ may be taken over unions of intervals; deduce from the proof of \cref{lem:graphon:weakreg} that $V(G)$ has a partition into $k\le4^{1/\eps^2}$ parts $V_i$ with $\bigl\lvert e(S,T)-\sum_{i,j}d(V_i,V_j)\abs{S\cap V_i}\abs{T\cap V_j}\bigr\rvert\le\eps n^2$ for all vertex sets $S,T$. (b) Deduce that the maximum number of edges across a cut of $G$ is within $\eps n^2$ of the maximum of $\sum_{i,j}d(V_i,V_j)\,s_i(\abs{V_j}-s_j)$ over integers $0\le s_i\le\abs{V_i}$, a problem in $k$ variables. Frieze and Kannan made a version of (a) algorithmic, and so approximated MAX-CUT within $\eps n^2$.

\item \consolidation\hard\label{ex:graphon:dense} \emph{Graphs are dense.} Prove part (iv) of \cref{thm:graphon:main}. (a) Approximate a graphon within $\eps$ in $L^1$ by a graphon $U$ that is constant on the squares of the partition of $[0,1]$ into $k$ equal intervals. (b) For $n=km$, let $G$ be the random graph on $\{1,\dots,n\}$ in which distinct $u$ and $v$ are adjacent, independently, with probability equal to the value of $U$ on $I_u\times I_v$; show that $\Pr[\graphoncut{W_G-U}>\eps+1/n]\to0$ as $m\to\infty$. \hint{For (a), use martingale convergence along dyadic partitions. For (b), take $S$ and $T$ to be vertex sets, use Hoeffding's inequality to bound by $2e^{-\eps^2n^2}$ the probability that $e(S,T)$ deviates from its mean by more than $\eps n^2$, and take a union bound over the $4^n$ pairs.}

\item \transfer\label{ex:graphon:words} \emph{Limits of words.} The picture of a word $s\in\{0,1\}^n$ is the function $f_s$ equal to $s_i$ on $[(i-1)/n,i/n)$. For measurable $f\colon[0,1]\to[0,1]$ and $w\in\{0,1\}^k$, let $t(w,f)$ be the probability that $k$ independent uniform points of $[0,1]$, sorted and coloured $1$ independently with probabilities $f(x)$, read $w$; for $f=f_s$ this samples $k$ positions of $s$. (a) Let $\norm{h}_I\defeq\sup_{a\le b}\abs{\int_a^bh}$. Show that uniformly random words $s_n$ of length $n$ have $\norm{f_{s_n}-\tfrac12}_1=\tfrac12$ but $\norm{f_{s_n}-\tfrac12}_I\to0$ in probability. (b) Show that $\abs{t(w,f)-t(w,g)}\le k^2\norm{f-g}_I$, and that the functions $f$, up to null sets, form a compact space under $\norm{f-g}_I$. (c) Conclude that every sequence of words has a subsequence along which all pattern densities converge, to those of some $f$. Why is no quotient by relabeling needed?

\item \transfer\hard\label{ex:graphon:permutons} \emph{Permutons.} A \term{permuton}\index{permuton} is a Borel probability measure on $[0,1]^2$ with uniform marginals; a permutation $\sigma$ of $\{1,\dots,n\}$ gives the permuton $\mu_\sigma$ with density $n$ on the squares $[(i-1)/n,i/n)\times[(\sigma(i)-1)/n,\sigma(i)/n)$. For a permutation $\tau$ of $\{1,\dots,k\}$, let $t(\tau,\mu)$ be the probability that $k$ independent points from $\mu$, listed by first coordinate, have second coordinates in the order $\tau$, and let $t(\tau,\sigma)$ be the fraction of $k$-sets of positions on which $\sigma$ has pattern $\tau$. (a) Show that $\abs{t(\tau,\mu_\sigma)-t(\tau,\sigma)}\le\binom k2/n$. (b) Show that $t(\tau,\cdot)$ is weakly continuous and that the permutons form a weakly compact set. (c) Conclude that every sequence of permutations with lengths tending to infinity has a subsequence along which all pattern densities converge to those of a permuton \parencite{HoppenEtAl2013}. Why does the weak topology suffice here, when graphs needed the cut norm? \hint{For (b), the events have boundaries of measure zero because the marginals have no atoms; use the portmanteau theorem and Prokhorov's theorem.}

\item \transfer\hard\label{ex:graphon:hypergraphs} \emph{What goes wrong for hypergraphs.} Encode $G(n,\tfrac12)$ by signs $\sigma_{uv}=-1$ for edges, $\sigma_{uv}=+1$ for non-edges, and $\sigma_{uu}=0$, and let $H$ be the $3$-uniform hypergraph of triples spanning an odd number of edges. (a) Show that every four vertices span an even number of edges of $H$, and that four given vertices span four with probability $\tfrac18$. (b) Show that for all vertex sets $X,Y,Z$, the number of triples in $X\times Y\times Z$ forming edges of $H$ is $\tfrac12\abs{X}\abs{Y}\abs{Z}+O(n^2\norm{\sigma})$, with $\norm{\sigma}$ the operator norm of the sign matrix. (c) Granting $\norm{\sigma}\le3\sqrt n$ with probability tending to $1$, conclude that $H$ is close to the constant $\tfrac12$ in the three-variable cut norm, while its expected density of complete $4$-vertex hypergraphs is $\tfrac18$, not $\tfrac1{16}$. Which extra variables would see the parity? \hint{For (b), fix $z$ and bound the sum over $x$ and $y$ of $\ind_X(x)\sigma_{xz}\,\sigma_{xy}\,\sigma_{yz}\ind_Y(y)$ by $n\norm{\sigma}$.}

\item \research\label{ex:graphon:gnn} \emph{Limits for learning.} (a) A graphon with a signal is a pair $(W,f)$ with $f\colon[0,1]\to[0,1]$; for a $1$-Lipschitz map $\rho\colon[0,1]\to[0,1]$, a layer of a mean-aggregation network maps it to $(W,\rho\circ T_Wf)$, where $T_Wf(x)=\int W(x,y)f(y)\dd y$. Show that $\norm{\rho\circ T_Wf-\rho\circ T_{W'}f'}_1\le2\graphoncut{W-W'}+\norm{f-f'}_1$. \textcite{Levie2023} builds a compact space of graphons with signals in which message-passing networks are Lipschitz, and derives generalization bounds from its covering numbers. Open-ended: what replaces this for sparse graphs, which all converge to the zero graphon? (b) Open-ended: in a network with two hidden layers of width $N$, the weights between the layers form an $N\times N$ matrix whose rows and columns can be relabeled independently. Find a metric, modeled on the cut distance, in which the network function is continuous and bounded sets are compact, as for the mean-field limit of one hidden layer (\cref{ch:transport}). Is training continuous in it?
\end{exercises}

\notesandrefs

The finite theory of \cref{sec:graphon:before} comes from \textcite{Mantel1907}, \textcite{Goodman1959}, \textcite{Szemeredi1978}, \textcite{RuzsaSzemeredi1978}, \textcite{Thomason1987}, \textcite{ChungGrahamWilson1989}, and \textcite{Gowers1997}; the limit theory from \textcite{FriezeKannan1999}, \textcite{LovaszSzegedy2006,LovaszSzegedy2007}, and \textcite{BorgsEtAl2008}. \textcite{Lovasz2012} is the standard monograph, which also covers sparse and bounded-degree limits, and \textcite{Zhao2023} is a textbook on regularity, quasirandomness, and graph limits whose compactness proof our sketch follows.

For flag algebras, start with \textcite{Razborov2007}; for the probabilistic side, with \textcite{DiaconisJanson2008}, who also give the references to Aldous and Hoover. The status of Sidorenko's conjecture \parencite{Sidorenko1993,OpenAI2026} should be checked afresh.

\printchapterbib
\end{refsection}
